\appendix
\originalchapter{来源与覆盖}\label{app:sources}
\originalsection{核心手册与章节对应}
以下页数是官方 PDF 的实际总页数, 含原有题名和 MIT OCW 许可末页; 它们不是中文正文页数. 原文件题名里的 Lecture 编号可能来自 2006 年等旧版本, 本册课次按 2019 年资源索引对应, 不擅自改写原文件日期. 中文讲义按数学主题重组, 下表表示内容来源与改编位置, 不表示逐字对照翻译.

\begin{longtable}{p{.8cm}p{8.2cm}p{.8cm}p{3cm}}
\toprule
编号 & 原手册题名 & 页数 & 本册对应\\
\midrule
\endhead
N01 & Floating-Point Arithmetic, the IEEE Standard (Per-Olof Persson) & 10 & 第 \ref{ch:float} 章\\
N02 & Some myths about floating-point arithmetic (Steven G. Johnson) & 2 & 第 \ref{ch:float} 章\\
N03 & Notes on the accuracy of naive summation (S. G. Johnson, 2019) & 3 & 第 \ref{ch:float} 章\\
N04 & Backwards stability of recursive summation (Steven G. Johnson, 2012/2015) & 3 & 第 \ref{ch:float} 章\\
N05 & Notes on the equivalence of norms (Steven G. Johnson, 2012) & 3 & 第 \ref{ch:float} 章\\
N06 & Gram-Schmidt Orthogonalization (Per-Olof Persson, 2006) & 3 & 第 \ref{ch:qr} 章\\
N07 & Householder Reflectors and Givens Rotations (Per-Olof Persson, 2006) & 13 & 第 \ref{ch:qr} 章\\
N08 & Hessenberg/Tridiagonal Reduction (Per-Olof Persson, 2006) & 2 & 第 \ref{ch:eigen} 章\\
N09 & The QR Algorithm I (Per-Olof Persson, 2006) & 4 & 第 \ref{ch:eigen} 章\\
N10 & The QR Algorithm II (Per-Olof Persson, 2006) & 9 & 第 \ref{ch:eigen} 章\\
N11 & Why restarting Arnoldi/Lanczos is not trivial (Steven G. Johnson, 2019) & 4 & 第 \ref{ch:krylov} 章\\
N12 & Sparse Matrix Algorithms (Per-Olof Persson, 2006) & 3 & 第 \ref{ch:sparse} 章\\
\midrule
合计 & 直接相关核心手册 & 59 & 12 份 PDF\\
\bottomrule
\end{longtable}
N01--N12 的官方链接在\href{https://ocw.mit.edu/courses/18-335j-introduction-to-numerical-methods-spring-2019/resources/lecture-notes/}{课程讲义资源栏}集中列出. 原件校验清单逐项记录了直链、文件名、页数、字节数与 SHA-256. 本册数学公式重新写成可编辑的原生 LaTeX, 没有用整页 PDF 截图代替正文.

\originalsection{性能资料与课堂主题}
第 \ref{ch:performance} 章另采用三份 Johnson 性能资料: Ideal-Cache Terminology (2 页), Performance Experiments with Matrix Multiplication (7 页), Experiments with Cache-Oblivious Matrix Multiplication (7 页), 共 16 页. 两份实验手册的硬件和编译器为历史环境, 本册保留其背景, 未把曲线当作新的实测结果.

本册的主课次映射为:
\begin{enumerate}
\item L2--4: 浮点、求和稳定性与范数, 对应第 \ref{ch:float} 章.
\item L5、L7: 条件数与扰动, 对应第 \ref{ch:condition} 章.
\item L7--8: SVD、回归、GSVD, 由 Alan Edelman 客座讲授, 对应第 \ref{ch:svd} 章.
\item L9--10: 正交化、Householder/Givens QR, 对应第 \ref{ch:qr} 章.
\item L11--12: 缓存与数据移动, 对应第 \ref{ch:performance} 章.
\item L13--14: LU、Cholesky、结构求解器, 对应第 \ref{ch:direct} 章.
\item L14--17: Schur、Hessenberg、幂法与 QR, 对应第 \ref{ch:eigen} 章.
\item L18--19: Arnoldi、Lanczos、Ritz 与重启, 对应第 \ref{ch:krylov} 章.
\item L20: GMRES 与收敛, 对应第 \ref{ch:gmres} 章.
\item L21--23: 预条件、CG、BiCG, 对应第 \ref{ch:cg} 章.
\item L24: 稀疏直接法, 对应第 \ref{ch:sparse} 章.
\end{enumerate}
上述主题均有官方逐周 HTML 说明, 但 L5、L7、L9、L13--14、L16--18、L20、L22--23 等课次没有覆盖整讲的独立手册. 一些课次列出补充阅读或往期 Persson 手册, 仍不等于完整逐讲讲稿. 本册因此补写了条件数定义与可达后向扰动、SVD 存在性与低秩逼近、LU/Cholesky 证明、Krylov 结构、GMRES 多项式论证、CG/PCG 推导和网格稀疏复杂度等内容. 这些证明与例题属于本册自编补充, 不标为 MIT 原文逐字译稿.

第 \ref{ch:experiment} 章的综合实验与 Julia 示例也是自编补充. 原课程 Julia 简介和速查表属于工具背景, 不是数值代数的数学证明来源. API 核对链接列入书目, 但未在编写环境运行 Julia, 因而没有声称其代码已做实机性能测试.

\originalsection{范围与许可边界}
官方讲义栏的 27 份 PDF 共 261 页, 涵盖整个 Introduction to Numerical Methods 的多个主题. 本册未纳入 L1 的求根、L6 的 ODE 以及 L25 以后优化、积分、Chebyshev 逼近和 FFT 的非数值代数全文, 以免与其他学科讲义重复. 原课提供的过往考试、问题集解答及外部版权教材也没有被混算为本册已经完整翻译的讲义体量.

主要改编材料依据 MIT OCW 默认 CC BY-NC-SA 4.0 许可, 仍需保留教师及手册作者、课程、题名、特殊署名、来源和改编说明. 课程页面明确提醒外部网站有独立条款, 不能因 MIT 链接了它们便默认继承 OCW 许可. 本册不下载或复制未获许可的 Trefethen--Bau 教材全文. Julia 简介中的 Project Jupyter 权利保留图像与优化概览中的 Springer 权利保留图像未进入本册改编或默认原件上传包.

本册的理论证明以所写明的精确算术和浮点模型为准. 成熟库的完整逐指令稳定性常数、Wilkinson 移位的专门全局收敛证明、多重网格与稳定化双共轭算法的完整实现, 只在明确条件和范围内作说明或引用, 不冒称已在这一本基础讲义中全部证明. 数值实验与机械排版检查也不替代数学证明.
